\section{Why the nanoscale is different}
\label{sec:fundamentals}
\begin{itemize}
  \item \textbf{Size} 1--100 nm; properties become discontinuous as the bulk contracts \slide{8}.
  \item \textbf{Quantum size effects}: discrete levels, band gap grows as size decreases (CdSe quantum dots: small = blue,
        large = red); melting point drops for small Au particles \slide{9--10}.
  \item \textbf{Surface effects}: $S/V\propto 1/\text{size}$ (cube of 1 m cut in 8 $\to$ area $\times2$, in 27 $\to$ $\times3$);
        55-atom cluster: 76\,\% surface atoms; Fe: half the atoms at the surface at $\sim$5 nm \slide{12--13}.
        Small, well-dispersed particles are therefore the best catalysts.
  \item \textbf{Goals of a synthesis} \slide{27}: monodisperse, same morphology, same composition and structure, no agglomeration;
        ideally fast, cheap, green, no critical raw materials.
\end{itemize}
